\chapter{Kiến trúc và Thiết kế RTL}
\label{ch:architecture}

Chương này trình bày cách các yêu cầu tại \cref{sec:req} được chuyển thành kiến trúc RTL. Nội dung tập
trung vào ranh giới khối, luồng transaction và các quyết định điều khiển chính. Bản đồ CSR, định dạng
Descriptor và bảng trạng thái đầy đủ được chuyển sang \cref{app:csr,app:desc,app:fsm}.

\section{Kiến trúc phân lớp}
\label{sec:arch-layers}

Thiết kế được chia thành ba lớp như \cref{fig:three-layer}. Lớp giao diện phần mềm tiếp nhận truy cập CSR,
lưu Command Descriptor và payload, đồng thời trả dữ liệu và Response Descriptor. Lớp Active Controller
giải mã lệnh, chọn Target, tạo chuỗi giao thức và xử lý lỗi. Lớp PHY đồng bộ tín hiệu vào và điều khiển
SDA/SCL. Cách phân lớp này giữ giao thức bus độc lập với cách phần mềm cấp lệnh.

\begin{figure}[H]
  \centering
  \begin{tikzpicture}[font=\small,node distance=0.35cm,
    layer/.style={draw,rounded corners,minimum width=13cm,minimum height=1.65cm,align=center}]
    \node[layer,fill=blue!8] (sw) {\textbf{Lớp giao diện phần mềm}\\CSR, DAT, Command/Response Queue, TX/RX Data Buffer};
    \node[layer,fill=green!9,below=of sw] (ctrl) {\textbf{Lớp Active Controller}\\\texttt{flow\_active}, ENTDAA, SCL generator, TX/RX serializer, bus monitor};
    \node[layer,fill=orange!10,below=of ctrl] (phy) {\textbf{Lớp PHY}\\đồng bộ đầu vào, output-enable và chuyển mạch Open Drain/Push-Pull};
    \draw[<->,thick] (sw)--(ctrl) node[midway,right]{Descriptor và handshake};
    \draw[<->,thick] (ctrl)--(phy) node[midway,right]{trạng thái bus và điều khiển pin};
  \end{tikzpicture}
  \caption{Kiến trúc ba lớp của I3C Active Controller.}
  \label{fig:three-layer}
\end{figure}

Thiết kế sử dụng một clock hệ thống và reset bất đồng bộ active-low. Các FSM, Queue và CSR cùng hoạt động
trong miền clock này. SCL được tạo bằng counter từ các giá trị timing trong CSR; do đó không hình thành
một clock domain nội bộ mới. SDA/SCL từ pad vẫn phải qua đồng bộ hai flip-flop trước khi được phân tích.

\section{Module cấp cao và kết nối khối}
\label{sec:arch-block}

Module \texttt{i3c\_controller\_top} là ranh giới IP. Phía Host gồm bus thanh ghi 32-bit. Phía bus gồm
\texttt{scl\_i/o}, \texttt{sda\_i/o}, output-enable của SDA và tín hiệu chọn chế độ. Bên trong, top-level
khởi tạo khối CSR, HCI Queue, Active Controller và PHY. Các module tương ứng là \texttt{csr\_registers},
\texttt{hci\_queues}, \texttt{controller\_active} và \texttt{i3c\_phy}. \Cref{fig:top-block} thể hiện các
kết nối chức năng thay vì liệt kê toàn bộ port.

\begin{figure}[H]
  \centering
  \resizebox{0.98\textwidth}{!}{\begin{tikzpicture}[>=Latex,font=\small,node distance=0.7cm and 0.9cm,
    b/.style={draw,rounded corners,minimum width=3.0cm,minimum height=1.35cm,align=center}]
    \node[b,fill=blue!8] (host) {Host / software\\register bus};
    \node[b,fill=blue!8,right=of host] (csr) {\texttt{csr\_registers}\\CSR + 32-entry DAT};
    \node[b,fill=blue!8,below=of csr] (q) {\texttt{hci\_queues}\\CMD, TX, RX, RESP};
    \node[b,fill=green!9,right=1.2cm of csr] (ca) {\texttt{controller\_active}\\protocol core};
    \node[b,fill=orange!10,right=of ca] (phy) {\texttt{i3c\_phy}\\2FF + pad control};
    \node[b,fill=gray!10,right=of phy] (bus) {I3C bus\\SCL, SDA};
    \draw[<->,thick] (host)--(csr);
    \draw[<->,thick] (csr)--(q);
    \draw[<->,thick] (q.east)--++(0.6,0)--++(0,1.35)--(ca.west);
    \draw[->,thick] (csr.east)--(ca.west) node[midway,above]{timing, DAT};
    \draw[<->,thick] (ca)--(phy);
    \draw[<->,thick] (phy)--(bus);
  \end{tikzpicture}}
  \caption{Sơ đồ khối cấp cao của \texttt{i3c\_controller\_top}.}
  \label{fig:top-block}
\end{figure}

Các Queue có chiều dữ liệu xác định. Command Queue và TX Data Buffer đi từ Host đến phần cứng. RX Data
Buffer và Response Queue đi theo chiều ngược lại. DAT có một cổng phần mềm và một cổng đọc phần cứng; cổng
đọc được phân kênh giữa \texttt{flow\_active} và \texttt{entdaa\_controller}. Sự phân kênh này tránh nhân
đôi bộ nhớ DAT và làm rõ quyền sở hữu ở từng pha giao dịch.

\section{Luồng transaction}
\label{sec:arch-flow}

\Cref{fig:dataflow} mô tả luồng nominal. Host cấu hình DAT và ghi Command Descriptor vào hai lần truy cập
32-bit. \texttt{csr\_registers} ghép hai DWORD rồi đẩy lệnh 64-bit vào Command Queue. Khi Controller được
enable, \texttt{flow\_active} lấy lệnh, yêu cầu đọc DAT và chọn nhánh I3C, I\textsuperscript{2}C, CCC hoặc
ENTDAA. Payload write được lấy từ TX Data Buffer; payload read được đóng gói theo DWORD vào RX Data Buffer.
Khi hoàn tất, \texttt{flow\_active} tạo Response Descriptor chứa error status, transaction ID và số byte.

\begin{figure}[H]
  \centering
  \begin{tikzpicture}[>=Latex,font=\small,node distance=0.65cm,
    f/.style={draw,rounded corners,minimum width=3.0cm,minimum height=0.9cm,align=center}]
    \node[f] (cmd) {Host ghi\\Command Descriptor};
    \node[f,right=of cmd] (cq) {Command Queue};
    \node[f,right=of cq] (fa) {\texttt{flow\_active}\\+ DAT lookup};
    \node[f,right=of fa] (wire) {SDA/SCL\\transaction};
    \node[f,below=of fa] (buf) {TX/RX Data Buffer};
    \node[f,left=of buf] (rq) {Response Queue};
    \node[f,left=of rq] (sw) {Host đọc\\data + response};
    \draw[->,thick] (cmd)--(cq); \draw[->,thick] (cq)--(fa); \draw[<->,thick] (fa)--(wire);
    \draw[<->,thick] (fa)--(buf); \draw[->,thick] (fa)--(rq); \draw[->,thick] (rq)--(sw);
    \draw[->,thick] (buf.west)-|(sw.south);
  \end{tikzpicture}
  \caption{Luồng dữ liệu từ Command Descriptor đến giao dịch bus và Response Descriptor.}
  \label{fig:dataflow}
\end{figure}

Handshake valid/ready được dùng tại ranh giới Queue. Producer giữ \texttt{valid} và dữ liệu ổn định cho
đến khi consumer chấp nhận. Quy tắc này quan trọng khi Response Queue đầy hoặc Host chưa đọc RX Data
Buffer. Đối với lệnh có \texttt{toc=0}, \texttt{flow\_active} có thể nhận Command Descriptor kế tiếp để tạo
chuỗi giao dịch liên tục. Nếu lệnh kế tiếp không có hoặc không được hỗ trợ, FSM kết thúc bus và trả response
thay vì chờ vô hạn.

\section{Clock, reset và đồng bộ tín hiệu}
\label{sec:arch-clk}

Clock hệ thống mục tiêu có chu kỳ 3~ns. Các tham số timing được lưu theo số chu kỳ, giúp cùng một cấu trúc
counter tạo được timing SDR và I\textsuperscript{2}C Fast Mode. Reset ngoài \texttt{rst\_ni} đưa FSM về
Idle và xóa trạng thái Queue. Software reset được tạo từ CSR, chỉ áp dụng cho các trạng thái điều khiển và
Queue được quy định; yêu cầu reset khi Controller đang bận được xử lý theo policy của CSR thay vì tác động
không đồng bộ đến bus.

\texttt{i3c\_phy} dùng hai flip-flop nối tiếp cho mỗi tín hiệu vào. Giá trị đã đồng bộ được đưa đến
\texttt{bus\_monitor}, \texttt{bus\_rx\_flow} và các FSM. Hai chu kỳ latency được chấp nhận vì clock hệ
thống nhanh hơn SCL. \Cref{fig:clk-sync} minh họa đường SDA; đường SCL có cấu trúc tương tự.

\begin{figure}[H]
  \centering
  \begin{tikzpicture}[>=Latex,font=\small,node distance=0.8cm,
    r/.style={draw,minimum width=1.5cm,minimum height=1.0cm,align=center}]
    \node[r] (pad) {SDA pad};
    \node[r,right=of pad] (ff1) {FF1\\metastability};
    \node[r,right=of ff1] (ff2) {FF2\\synchronized};
    \node[r,right=of ff2] (mon) {bus monitor\\RX flow};
    \draw[->,thick] (pad)--(ff1); \draw[->,thick] (ff1)--(ff2); \draw[->,thick] (ff2)--(mon);
    \draw[dashed,->] ([yshift=-0.75cm]ff1.south)--++(0,-0.6) node[below]{\texttt{clk\_i}};
    \draw[dashed,->] ([yshift=-0.75cm]ff2.south)--++(0,-0.6) node[below]{\texttt{clk\_i}};
  \end{tikzpicture}
  \caption{Đồng bộ hai flip-flop cho tín hiệu bus đi vào miền clock hệ thống.}
  \label{fig:clk-sync}
\end{figure}

\section{Ranh giới với thiết kế tham chiếu}
\label{sec:arch-reference}

CHIPS Alliance \texttt{i3c-core} hỗ trợ bối cảnh tích hợp rộng, gồm các adapter hệ thống, cấu trúc HCI và
cả những nhánh chức năng ngoài phạm vi khóa luận~\cite{chipsalliance_i3c}. Thiết kế trong khóa luận không
nhằm thay thế một implementation đầy đủ. Phần kế thừa về ý tưởng gồm cách tách lớp PHY/Controller, mô hình
Descriptor và một số FSM truyền nhận. Phần được thu hẹp gồm bus Host đơn giản, DAT 32 mục, bốn Queue đồng
bộ và một Active Controller chỉ hỗ trợ tập tính năng tại \cref{tab:req-func}.

Thay đổi quan trọng nhất nằm ở đường phát lệnh. FSM trung tâm được viết lại để xử lý đầy đủ Command
Descriptor mục tiêu, Immediate Data, CCC, ENTDAA, I\textsuperscript{2}C, response và abort. Vì repository
không chứa một snapshot đã pin của toàn bộ mã nguồn tham chiếu, so sánh định lượng theo số dòng không được
suy ra từ tài liệu mô tả. Việc đánh giá tại \cref{ch:results} phải dùng baseline có thể tái tạo và tách
riêng nhận xét định tính.

\section{Queue, DAT và Descriptor}
\label{sec:arch-desc}

\texttt{hci\_queues} khởi tạo bốn \texttt{sync\_fifo}. Độ rộng Command Queue là 64 bit, còn TX Data
Buffer, RX Data Buffer và Response Queue là 32 bit. Độ sâu mặc định ở top-level là 64; Testbench giảm
độ sâu xuống 8 để kích hoạt nhanh trạng thái full/empty. Mỗi \texttt{sync\_fifo} dùng read pointer và
write pointer có thêm một bit wrap. Elaborating với độ sâu không phải lũy thừa hai bị chặn bằng immediate
assertion.

DAT có 32 mục. Mỗi mục cung cấp Dynamic Address hoặc Static Address và cờ xác định Target
I\textsuperscript{2}C. \texttt{flow\_active} dùng \texttt{dev\_idx} trong Descriptor để chọn mục; ENTDAA
dùng chỉ số bắt đầu và số Target từ Address Assignment Descriptor. Các bitfield đầy đủ được trình bày tại
\cref{app:desc}; \cref{tab:desc-summary} chỉ giữ thông tin cần cho luồng kiến trúc.

\begin{table}[H]
  \centering
  \caption{Vai trò của các cấu trúc giao tiếp chính.}
  \label{tab:desc-summary}
  \begin{tabular}{L{3.8cm} C{2.0cm} L{7.3cm}}
    \toprule
    \textbf{Cấu trúc} & \textbf{Độ rộng} & \textbf{Nội dung sử dụng} \\
    \midrule
    Command Descriptor & 64 bit & Loại lệnh, hướng, chế độ, \texttt{dev\_idx}, transaction ID, độ dài và cờ kết thúc. \\
    Response Descriptor & 32 bit & Error status, transaction ID và số byte đã xử lý. \\
    DAT entry & 32 bit & Static/Dynamic Address và thuộc tính loại Target. \\
    TX/RX Data Buffer word & 32 bit & Tối đa bốn byte payload theo thứ tự little-endian trong word. \\
    \bottomrule
  \end{tabular}
\end{table}

Command Descriptor chấp nhận Regular Transfer, Immediate Data Transfer và Address Assignment. Combo
Transfer và mode khác \texttt{sdr0} bị trả \texttt{NotSupported}. Immediate Data chỉ cho phép write và
tối đa bốn byte. Address Assignment phải mang ENTDAA, có \texttt{toc=1},
\texttt{wroc=1}, số Target khác không và không vượt quá biên DAT.

\section{Mô hình xử lý lỗi}
\label{sec:arch-err}

Lỗi được tích lũy trong suốt transaction và ánh xạ sang trường \texttt{ERR\_STATUS} khi ghi Response
Descriptor. Đây là convention của interface TCRI/HCI trong implementation, không phải mã lỗi do MIPI I3C
Basic định nghĩa. Khi nhiều điều kiện cùng xuất hiện, hàm \texttt{map\_resp\_err\_status()} áp dụng thứ tự
ưu tiên cố định; vì vậy response chỉ mang một mã.

\begin{table}[H]
  \centering
  \caption{Các error status được RTL tạo ra.}
  \label{tab:err-status}
  \begin{tabular}{C{1.5cm} L{4.0cm} L{7.5cm}}
    \toprule
    \textbf{Mã} & \textbf{Tên RTL} & \textbf{Điều kiện chính} \\
    \midrule
    0x0 & \texttt{Success} & Transaction kết thúc không có lỗi được ghi nhận. \\
    0x4 & \texttt{AddrHeader} & Target NACK Address Header. \\
    0x5 & \texttt{Nack} & Dynamic Address bị từ chối hai lần trong ENTDAA. \\
    0x6 & \texttt{Ovl} & TX underflow hoặc RX overflow. \\
    0x7 & \texttt{I3cShortReadErr} & Target kết thúc read sớm khi policy yêu cầu báo lỗi. \\
    0x8 & \texttt{HcAborted} & Host yêu cầu abort transaction. \\
    0x9 & \shortstack[l]{\texttt{I2cDataNackOr}\\\texttt{I3cBusAborted}} & Data NACK của I\textsuperscript{2}C hoặc bus abort tương ứng. \\
    0xA & \texttt{NotSupported} & Descriptor hoặc mode không thuộc tập hỗ trợ. \\
    \bottomrule
  \end{tabular}
\end{table}

Address NACK có ưu tiên cao nhất, tiếp theo là data NACK, lỗi ENTDAA, overflow/underflow, short read,
unsupported và abort. Cách ưu tiên này bảo đảm lỗi gần giao thức nhất không bị che bởi yêu cầu abort đến
muộn. Response Queue sử dụng valid-before-ready: \texttt{WriteResp} giữ \texttt{valid} và dữ liệu cho đến
khi Queue chấp nhận.

\section{Các khối RTL hỗ trợ}

\subsection{PHY và bus monitor}
\label{sec:rtl-phy}
\label{sec:rtl-monitor}

\texttt{i3c\_phy} thực hiện hai nhiệm vụ: đồng bộ SDA/SCL vào clock hệ thống và chuyển tín hiệu điều khiển
nội bộ thành output/output-enable. Với Open Drain, mức cao được tạo bằng việc nhả SDA. Với Push-Pull,
output-enable được giữ và giá trị logic được điều khiển trực tiếp. PHY không tự quyết định pha giao thức;
quyền chọn mode thuộc khối tích hợp Active Controller.

\texttt{bus\_monitor} nhận tín hiệu đã đồng bộ, tạo pulse cạnh và phát hiện START, Repeated START, STOP.
Hai module phụ \texttt{edge\_detector} và \texttt{stable\_high\_detector} lọc thời điểm đánh giá. Monitor
chỉ báo sự kiện; nó không thay đổi SDA/SCL. Việc tách phát hiện khỏi tạo tín hiệu giúp Scoreboard và SVA có
thể quan sát cùng định nghĩa sự kiện với RTL.

\subsection{CSR, DAT và HCI queues}
\label{sec:rtl-csr}
\label{sec:rtl-queues}

\texttt{csr\_registers} hiện thực 26 địa chỉ CSR cùng vùng DAT 32 mục. Khối tạo default timing, giữ bit
enable/abort, phản ánh trạng thái bận và chuyển truy cập tại các data port thành handshake Queue. Command
Descriptor 64-bit được ghép từ hai lần ghi 32-bit. Nếu có backpressure, dữ liệu và \texttt{valid} được giữ
đến khi \texttt{hci\_queues} chấp nhận. DAT write áp dụng mask để các bit reserved luôn bằng không.

\texttt{hci\_queues} không diễn giải nội dung. Bốn FIFO đồng bộ dùng chung clock và hỗ trợ flush bằng
software reset. Ngoài full/empty, wrapper xuất depth hiện tại để CSR tạo trạng thái Queue. Thiết kế này đủ
cho PIO flow trong phạm vi, nhưng không cung cấp threshold interrupt hoặc DMA như một HCI đầy đủ.

\subsection{SCL generator}
\label{sec:rtl-scl}

\texttt{scl\_generator} có 14 trạng thái, được chia thành bốn nhóm chức năng. Nhóm START tạo cạnh SDA và
giữ setup/hold; nhóm clock luân phiên mức thấp, mức cao và chờ lệnh; nhóm Repeated START tạo lại điều kiện
START khi transaction chưa có STOP; nhóm STOP đưa SDA lên cao và chờ Bus Free Condition. Counter được nạp
từ CSR theo trạng thái hiện tại. Tên và chuyển tiếp đầy đủ được liệt kê tại \cref{app:fsm}.

\begin{figure}[H]
  \centering
  \begin{tikzpicture}[>=Latex,font=\scriptsize,node distance=0.45cm and 0.4cm,
    s/.style={draw,rounded corners,minimum width=2.3cm,minimum height=0.65cm,align=center}]
    \node[s] (idle) {Idle};
    \node[s,right=of idle] (start) {GenerateStart\\SdaFall, HoldStart};
    \node[s,right=of start] (clk) {DriveLow\\DriveHigh, WaitCmd};
    \node[s,below=of clk] (rs) {GenerateRstart\\SclHigh, SdaFall};
    \node[s,left=of rs] (stop) {GenerateStop\\SclHigh, SdaRise};
    \node[s,left=of stop] (free) {BusFree};
    \draw[->] (idle)--(start); \draw[->] (start)--(clk);
    \draw[->] (clk)--(rs); \draw[->] (rs)--(clk);
    \draw[->] (clk)--(stop); \draw[->] (stop)--(free); \draw[->] (free)--(idle);
  \end{tikzpicture}
  \caption{Nhóm trạng thái chức năng của FSM \texttt{scl\_generator}.}
  \label{fig:scl-fsm}
\end{figure}

Tín hiệu \texttt{takeover\_i} tạo fast path từ \texttt{DriveHigh} đến \texttt{RstartSdaFall}. Đường này
được dùng khi Active Controller phải giành lại SDA tại ranh giới T-Bit hoặc ENTDAA mà không chờ chuỗi
Repeated START thông thường. \texttt{scl\_use\_od\_low\_i} chọn timing low dành cho pha Open Drain; việc
chọn không dựa riêng vào mode I3C/I\textsuperscript{2}C mà dựa vào pha thực tế.

\subsection{TX/RX serializer}
\label{sec:rtl-txrx}

\texttt{bus\_tx\_flow} nhận yêu cầu phát một bit hoặc một byte. FSM bốn trạng thái phân xử yêu cầu, kích
hoạt \texttt{bus\_tx} và trả pulse hoàn tất. \texttt{bus\_tx} có năm trạng thái để chờ cạnh SCL, đặt SDA,
giữ dữ liệu và chuyển bit. Output \texttt{sda\_drive\_o} được tách khỏi giá trị SDA để
\texttt{controller\_active} tạo output-enable đúng ở cả Open Drain và Push-Pull.

\texttt{bus\_rx\_flow} cũng có bốn trạng thái. Khối chốt bit tại cạnh phù hợp, dịch theo thứ tự MSB-first
và trả byte hoàn chỉnh. Yêu cầu bit handoff là biến thể của đọc một bit: sau khi lấy mẫu, quyền điều khiển
SDA có thể được chuyển ngay cho Controller. Cơ chế này được dùng cho ACK và T-Bit, nơi một chu kỳ trễ có
thể làm sai pha.

\begin{figure}[H]
  \centering
  \begin{tikzpicture}[font=\small,node distance=0.55cm and 0.65cm,
    s/.style={draw,rounded corners,minimum width=2.4cm,minimum height=0.72cm,align=center}]
    \node[s] (txf) {\texttt{bus\_tx\_flow}\\4 states};
    \node[s,right=of txf] (tx) {\texttt{bus\_tx}\\5 states};
    \node[s,right=of tx] (sda) {SDA output\\+ drive enable};
    \node[s,below=of tx] (rx) {\texttt{bus\_rx\_flow}\\4 states};
    \draw[->,thick] (txf)--(tx); \draw[->,thick] (tx)--(sda);
    \draw[<-,thick] (rx.east)-|(sda.south); \draw[->,thick] (rx)--(txf) node[midway,left]{done/data};
  \end{tikzpicture}
  \caption{Quan hệ giữa các FSM truyền và nhận nối tiếp.}
  \label{fig:txrx-fsm}
\end{figure}

\section{Phân hệ ENTDAA}
\label{sec:rtl-entdaa}

ENTDAA được chia thành hai FSM. \texttt{entdaa\_controller} có 7 trạng thái và quản lý vòng lặp qua các
mục DAT. \texttt{entdaa\_fsm} có 8 trạng thái và xử lý một Target: phát reserved address ở hướng Read,
nhận ACK, thu 64 bit PID/BCR/DCR, phát Dynamic Address cùng parity và nhận ACK. Việc chia hai cấp làm rõ
khác biệt giữa vòng lặp phần mềm yêu cầu và giao thức của một vòng arbitration.

\begin{figure}[H]
  \centering
  \begin{tikzpicture}[>=Latex,font=\small,node distance=0.55cm and 0.8cm,
    s/.style={draw,rounded corners,minimum width=3.0cm,minimum height=0.85cm,align=center}]
    \node[s,fill=green!8] (loop) {\texttt{entdaa\_controller}\\7-state loop manager};
    \node[s,fill=green!8,right=of loop] (round) {\texttt{entdaa\_fsm}\\8-state per-Target round};
    \node[s,above=of loop] (dat) {DAT\\address candidate};
    \node[s,below=of loop] (flow) {\texttt{flow\_active}\\Descriptor + result write};
    \node[s,right=of round] (bus) {TX/RX flow\\SDA/SCL};
    \draw[<->,thick] (loop)--(round); \draw[<->,thick] (loop)--(dat);
    \draw[<->,thick] (loop)--(flow); \draw[<->,thick] (round)--(bus);
  \end{tikzpicture}
  \caption{Phân cấp điều khiển của phân hệ ENTDAA.}
  \label{fig:entdaa-fsm}
\end{figure}

\texttt{entdaa\_controller} bắt đầu tại chỉ số \texttt{dev\_idx}, đọc Dynamic Address ứng viên từ DAT và
giới hạn tổng chỉ số trong 32 mục. Nếu Target NACK Dynamic Address lần đầu, FSM thử lại cùng vòng. NACK lần
hai kết thúc với mã lỗi 0x5. Nếu không có Target phản hồi reserved address, quá trình kết thúc bình thường
với số lượng kết quả đã thu được. Mỗi kết quả được \texttt{flow\_active} ghi thành ba DWORD trong RX Data
Buffer: PID[47:16], PID[15:0]/BCR/DCR và Dynamic Address.

\begin{figure}[H]
  \centering
  \begin{tikzpicture}[>=Latex,font=\small,node distance=0.45cm,
    p/.style={draw,rounded corners,minimum width=9.4cm,minimum height=0.78cm,align=center}]
    \node[p] (a) {Đọc Dynamic Address ứng viên từ DAT};
    \node[p,below=of a] (b) {Phát Repeated START và chạy một vòng \texttt{entdaa\_fsm}};
    \node[p,below=of b] (c) {Nhận PID/BCR/DCR, phát Dynamic Address và lấy ACK};
    \node[p,below=of c] (d) {Ghi ba DWORD kết quả vào RX Data Buffer};
    \node[p,below=of d] (e) {Tăng chỉ số; lặp nếu còn Target và chưa đủ \texttt{dev\_count}};
    \draw[->,thick] (a)--(b); \draw[->,thick] (b)--(c); \draw[->,thick] (c)--(d); \draw[->,thick] (d)--(e);
    \draw[->] (e.east)--++(1.0,0)|-(a.east);
  \end{tikzpicture}
  \caption{Luồng xử lý một tập Target trong ENTDAA.}
  \label{fig:entdaa-loop}
\end{figure}

\section{FSM trung tâm \texttt{flow\_active}}
\label{sec:rtl-flow-active}

\subsection{Vai trò và tổ chức trạng thái}
\label{sec:flow-states}

\texttt{flow\_active} là khối điều phối chính. FSM có 13 trạng thái, nhưng chi tiết byte/bit không được
mã hóa thành thêm state. Biến \texttt{issue\_phase\_q} làm micro-counter cho các pha START, Address Header,
ACK, data và T-Bit. Cách tổ chức này giữ state count cố định khi thêm biến thể transaction, đồng thời cho
phép chia sẻ \texttt{IssueCmd} giữa I3C, I\textsuperscript{2}C, read, write và Immediate Data.

\begin{table}[H]
  \centering
  \caption{Vai trò tóm tắt của 13 trạng thái \texttt{flow\_active}.}
  \label{tab:flow-states}
  \begin{tabular}{C{1.0cm} L{3.5cm} L{8.4cm}}
    \toprule
    \textbf{ID} & \textbf{State} & \textbf{Vai trò} \\
    \midrule
    0 & \texttt{Idle} & Xóa context transaction và giữ bus ở trạng thái nghỉ. \\
    1 & \texttt{WaitForCmd} & Chờ Command Queue có Descriptor hợp lệ. \\
    2 & \texttt{FetchDAT} & Kiểm tra Descriptor, phát yêu cầu đọc DAT. \\
    3 & \texttt{WaitDAT} & Chờ dữ liệu DAT và quyết định nhánh giao thức. \\
    4 & \texttt{I3CBcastHeader} & Phát broadcast preamble khi cấu hình yêu cầu. \\
    5 & \shortstack[l]{\texttt{IssueImmediate}\\\texttt{Ccc}} & Phát Broadcast/Direct ENEC hoặc DISEC. \\
    6 & \texttt{FetchTxData} & Lấy DWORD tiếp theo từ TX Data Buffer. \\
    7--10 & \texttt{Init*} & Phát Address Header cho I3C/I\textsuperscript{2}C read/write. \\
    11 & \texttt{IssueCmd} & Truyền payload, chạy ENTDAA và xử lý continuation/abort. \\
    12 & \texttt{WriteResp} & Ghi Response Descriptor bằng valid/ready handshake. \\
    \bottomrule
  \end{tabular}
\end{table}

Bảng chuyển trạng thái đầy đủ được đưa vào \cref{app:fsm}. \Cref{fig:flow-fsm} giữ các nhánh chính để thể
hiện câu chuyện thiết kế: chuẩn bị Descriptor, chọn preamble, thực hiện transaction và trả response.

\begin{figure}[H]
  \centering
  \resizebox{0.98\textwidth}{!}{\begin{tikzpicture}[>=Latex,font=\scriptsize,node distance=0.45cm and 0.55cm,
    s/.style={draw,rounded corners,minimum width=2.35cm,minimum height=0.68cm,align=center}]
    \node[s] (idle) {Idle};
    \node[s,right=of idle] (wait) {WaitForCmd};
    \node[s,right=of wait] (fetch) {FetchDAT};
    \node[s,right=of fetch] (wd) {WaitDAT};
    \node[s,above right=of wd] (bcast) {I3CBcastHeader};
    \node[s,right=of bcast] (ccc) {IssueImmediateCcc};
    \node[s,below right=of wd] (i3cw) {InitI3CWrite};
    \node[s,right=of i3cw] (i3cr) {InitI3CRead};
    \node[s,below=of i3cw] (i2cw) {InitI2CWrite};
    \node[s,right=of i2cw] (i2cr) {InitI2CRead};
    \node[s,right=1.0cm of i3cr] (tx) {FetchTxData};
    \node[s,right=of tx] (issue) {IssueCmd};
    \node[s,right=of issue] (resp) {WriteResp};
    \draw[->] (idle)--(wait); \draw[->] (wait)--(fetch); \draw[->] (fetch)--(wd);
    \draw[->] (wd)--(bcast); \draw[->] (bcast)--(ccc); \draw[->] (bcast)--(i3cw);
    \draw[->] (wd)--(i3cw); \draw[->] (wd)--(i3cr); \draw[->] (wd)--(i2cw); \draw[->] (wd)--(i2cr);
    \draw[->] (i3cw)--(tx); \draw[->] (i2cw)--(tx); \draw[->] (tx)--(issue);
    \draw[->] (i3cr)--(issue); \draw[->] (i2cr)--(issue); \draw[->] (ccc.east)-|([yshift=0.2cm]resp.north);
    \draw[->] (issue)--(resp); \draw[->] (resp.south) |- ++(0,-1.0) -| (idle.south);
    \draw[->,dashed] (issue.north) to[bend left=18] node[above]{continuation} (fetch.north);
    \draw[->,dashed] (fetch.south) to[bend right=15] node[below]{invalid} (resp.south);
  \end{tikzpicture}}
  \caption{Các chuyển tiếp chính của FSM \texttt{flow\_active} 13 trạng thái.}
  \label{fig:flow-fsm}
\end{figure}

\subsection{Giải mã và định tuyến Command Descriptor}
\label{sec:flow-route}

Tại \texttt{FetchDAT}, ba phép kiểm tra xảy ra trước khi tác động bus: loại Descriptor có được hỗ trợ hay
không, mode có phải \texttt{sdr0} hay không, và các trường riêng có hợp lệ hay không. Descriptor lỗi đi
thẳng đến \texttt{WriteResp} với \texttt{NotSupported}. Descriptor hợp lệ phát yêu cầu đọc DAT; state
\texttt{WaitDAT} tách một chu kỳ để dữ liệu ổn định.

Sau DAT lookup, cờ loại Target quyết định timing và cơ chế ACK. Nếu bật broadcast header, private transfer
I3C đi qua \texttt{I3CBcastHeader}; nếu không, FSM vào trực tiếp state \texttt{InitI3C*}. Immediate CCC
luôn dùng broadcast header. Address Assignment phát CCC ENTDAA rồi trao quyền vòng lặp cho phân hệ ENTDAA.

\begin{figure}[H]
  \centering
  \begin{tikzpicture}[>=Latex,font=\small,node distance=0.42cm,
    p/.style={draw,rounded corners,minimum width=9.2cm,minimum height=0.75cm,align=center}]
    \node[p] (a) {Pop Command Descriptor và kiểm tra \texttt{attr}/\texttt{mode}/field};
    \node[p,below=of a] (b) {Đọc DAT tại \texttt{dev\_idx}};
    \node[p,below=of b] (c) {Chọn I3C hoặc I\textsuperscript{2}C; chọn private, CCC hoặc ENTDAA};
    \node[p,below=of c] (d) {Phát preamble và Address Header};
    \node[p,below=of d] (e) {Chạy payload trong \texttt{IssueCmd}};
    \node[p,below=of e] (f) {Tạo STOP/continuation và ghi Response Descriptor};
    \draw[->,thick] (a)--(b); \draw[->,thick] (b)--(c); \draw[->,thick] (c)--(d);
    \draw[->,thick] (d)--(e); \draw[->,thick] (e)--(f);
  \end{tikzpicture}
  \caption{Luồng phát lệnh của \texttt{flow\_active}.}
  \label{fig:flow-issue}
\end{figure}

\subsection{Private write, private read và Immediate Data}
\label{sec:flow-merge}

Regular write lấy từng DWORD từ TX Data Buffer. \texttt{tx\_byte\_idx\_q} chọn byte hiện tại; sau mỗi
byte I3C, FSM phát odd-parity T-Bit. Với I\textsuperscript{2}C, cùng state phát byte nhưng yêu cầu
\texttt{bus\_rx\_flow} lấy ACK/NACK. Khi DWORD cuối không đủ bốn byte, \texttt{remaining\_len\_q} dừng
đường truyền đúng tại byte cuối.

Private read nhận byte vào \texttt{rx\_dword\_q}. Dữ liệu được commit khi đủ bốn byte, hết độ dài, Target
kết thúc sớm hoặc abort. Với I3C, bit sau dữ liệu được lấy qua handoff path; với I\textsuperscript{2}C,
Controller phát ACK khi còn dữ liệu và NACK ở byte cuối. RX full tạo \texttt{Ovl} và chấm dứt transaction
theo đúng pha thay vì ghi đè dữ liệu chưa đọc.

Immediate Data không dùng TX Data Buffer. Các byte được lấy trực tiếp từ DWORD thứ hai của Descriptor,
nhưng vẫn đi qua \texttt{IssueCmd}; do đó parity, Address Header, response length và abort dùng chung logic
với Regular write. CCC immediate được tách sang \texttt{IssueImmediateCcc} vì Direct CCC cần chuỗi CCC,
Repeated START, Target Address và event byte khác private transfer.

\subsection{Continuation, abort và kết thúc transaction}
\label{sec:flow-end}

\texttt{toc} quy định việc kết thúc sau Command Descriptor. Với \texttt{toc=1}, FSM tạo STOP và trả
response khi \texttt{wroc=1} hoặc khi có lỗi. Với \texttt{toc=0}, FSM kiểm tra Descriptor kế tiếp. Chỉ
Regular SDR private transfer hợp lệ mới được nối tiếp; trường hợp khác tạo STOP và báo lỗi để tránh giữ bus
không xác định.

Abort không được xử lý bằng cách nhảy tức thời về Idle. Nếu transaction đang ở pha cần kết thúc có trật tự,
FSM chờ TX/RX rảnh, tạo STOP và sau đó trả \texttt{HcAborted}. Với I3C read, Controller phải hoàn tất việc
lấy T-Bit và takeover trước khi kéo bus; với I\textsuperscript{2}C read, Controller phát NACK rồi STOP.
Cách xử lý theo pha ngăn xung SDA bất hợp lệ trong khi SCL cao.

\section{Tích hợp trong \texttt{controller\_active}}
\label{sec:rtl-controller-active}

\texttt{controller\_active} tích hợp FSM phát lệnh, SCL generator, các TX/RX flow, bus monitor và phân hệ
ENTDAA. Các yêu cầu bus từ \texttt{flow\_active} và ENTDAA dùng chung
serializer qua MUX. DAT cũng được phân kênh theo \texttt{daa\_active}. Một latch giữ yêu cầu Repeated START
ENTDAA cho đến khi SCL generator xác nhận, tránh mất pulse khi hai FSM không hoàn tất cùng chu kỳ.

MUX SDA có thứ tự ưu tiên rõ ràng. Khi SCL generator đang tạo START, Repeated START hoặc STOP, giá trị SDA
của generator được chọn. Trong pha dữ liệu, serializer hoặc ENTDAA điều khiển SDA. Tín hiệu mode và
\texttt{sda\_drive} quyết định output-enable: Open Drain chỉ enable khi phát mức thấp; Push-Pull enable ở
cả hai mức. \Cref{fig:od-pp-switch} mô tả logic ở mức chức năng.

\begin{figure}[H]
  \centering
  \begin{tikzpicture}[>=Latex,font=\small,node distance=0.65cm,
    b/.style={draw,rounded corners,minimum width=3.2cm,minimum height=0.9cm,align=center}]
    \node[b] (src) {SDA source MUX\\SCL generator / TX / ENTDAA};
    \node[b,right=of src] (mode) {Mode decision\\Open Drain / Push-Pull};
    \node[b,right=of mode] (oe) {\texttt{sda\_o} + \texttt{sda\_oe}};
    \node[b,below=of mode] (rule) {OD: enable khi SDA=0\\PP: enable trong toàn pha};
    \draw[->,thick] (src)--(mode); \draw[->,thick] (mode)--(oe); \draw[->,thick] (rule)--(mode);
  \end{tikzpicture}
  \caption{Chuyển mạch nguồn SDA và điều khiển Open Drain/Push-Pull.}
  \label{fig:od-pp-switch}
\end{figure}

Kiến trúc tích hợp bảo đảm chỉ một nguồn điều khiển SDA trong mỗi pha. Các ràng buộc mutual exclusion và
handoff được kiểm tra bằng SVA ở \cref{sec:ver-sva}; hành vi xuyên suốt được kiểm tra bởi Scoreboard ở
\cref{sec:ver-tlm}.
